%!TEX root = Physics231.tex
\chapter{Spring Theory}
\label{ch:spring:theory}

\begin{quote}
\emph{Everything is a harmonic oscillator.}
\\ -- every physics student, eventually.\sidenote{At least approximately.}
\end{quote}


\section{Newton}

You know from your earliest physics memories that a harmonic oscillator is a system with a canonical kinetic term and a quadratic potential. The first time you saw it may have been in the context of, say, a spring\sidenote{Or a pendulum.} with some spring constant $k$. This means that the potential is $V(x)=(k/2)x^2$ where $x$ is the displacement from the equilibrium configuration that minimizes potential energy. 

The variable $x$ describes the state of a system. This state evolves in time, so it is a function $x(t)$. Physics tells us how this state evolves in time.\sidenote{This is too arrogant. What I really mean by `physics' is that our mathematical description of physics. Perhaps there is a different mathematical description wherein the evolution rule is not infinitesimal.} It does this by telling us how the state $x(t_0)$ at some time changes with an \emph{infinitesimal} change in time, $t_0 \to t_0 + \D{t}$. This is called the equation of motion. 

In Newtonian physics, the equation of motion is $F=m\ddot{x}$, where $F = -V'(x)$. This gives
\begin{align}
    \frac{ \D[2]{x(t)} }{ \D{t}^2 } 
    = 
    - \frac{ k }{ m } x(t) \ .
\end{align}
The rule for the infinitesimal change is \emph{second order}. Those who are well versed in differential equations will recall that this means you need first order initial data, like the initial velocity.\sidenote{The fancy way of saying this is Picard's theorem for the existence and uniqueness of ordinary differential equations. When I was an undergrad, I took the math-for-mathematicians first year calculus sequence. While my friends taking engineering math learned how to solve a handful of common differential equations, we spent 10 weeks learning Picard's construction to argue that solutions exist. If you wonder why I often make jokes about mathematicians and engineers, this experience is probably why.}


\section{Lagrange, traditionally}
\label{sec:lagrange:spring:theory}

You meet the harmonic oscillator again when upgrading from Newtonian to Lagrangian dynamics.\sidenote{For our purposes, we treat Lagrangian and Hamiltonian dynamics as the same up to a Legendre transform, see Appendix~\ref{ch:Lagrange:Legendre}.} 
% 
One of the fundamental principles in physics is the \textbf{principle of least action}\index{principle of least action}. This is the statement that classically, the trajectory of a system is governed by the minimum of a functional called the action, $S$. A functional is a ``function of functions.'' The position of a particle $\vec{x}(t)$ is a function of time. The action is a function of $\vec{x}(t)$. You may be more familiar with the action as the time integral of the Lagrangian,
\begin{align}
    S = \int \D{t}\, L[\vec{x}, \dot{\vec{x}}] \ ,
\end{align}
where we use the standard convention where we treat the position and velocity as independent variables in phase space. In quantum mechanics, this same object appears as a weight for possible trajectories.\footnote{I recommend looking up Feynman's phasor description of the sum over histories to see how the action in quantum mechanics leads to the principle of least action in classical mechanics.}


In analytical mechanics we derive the equation of motion by varying the Lagrangian. This is called the \textbf{variational priciple}\index{variational principle}. The standard treatment goes something like this. The Lagrangian of a classical particle is the kinetic minus the potential energy,
\begin{align}
    L = \frac{m}{2} \dot{\vec{x}}^2 - V[\vec{x}] \ .
\end{align}
When we treat $\vec{x}$ and $\dot{\vec{x}}$ as independent variables, the infinitesimal change (variation) of the action $S$ with with respect to infinitesimal perturbations $\delta\vec{x}$ and $\delta\dot{\vec{x}}$ is
\begin{align}
    \delta S =
    \int \D{t} \,
    \left[
    \frac{\partial L}{\partial \vec{x}} \delta\vec{x}
    +
    \frac{\partial L}{\partial \dot{\vec{x}}}\delta\dot{\vec{x}}
    \right] \ .
\end{align}
Now we `remember' that $\dot{\vec{x}}$ is actually the time derivative of $\vec{x}(t)$. This means we can integrate the second term by parts to remove the time derivative from $\delta\dot{\vec{x}}$. This is simply the statement that there is \emph{one} variation we're doing: the function $\delta\vec{x}(t)$. So if the path changes by $\vec{x}\to\vec{x}+\delta\vec{x}$, then the time derivative of the path changes by
\begin{align}
    \dot{\vec{x}} 
    \to \dot{\vec{x}} 
    + \frac{\D{}}{\D{t}}(\delta\vec{x}) 
    \equiv \dot{\vec{x}} + \delta\dot{\vec{x}} \ . 
\end{align}
% 
% 
We assume some boundary conditions---the position of the particle at some initial and some final times. This means we can integrate by parts because the boundary terms in the integration by parts vanishes.  We are thus left with (suppressing the limits of integration)
\begin{align}
    \delta S = \int dt\, 
    \left[
    \frac{\partial L}{\partial \vec{x}}
    - 
    \frac{\D{}}{\D{t}}
    \frac{\partial L}{\partial \dot{\vec{x}}} 
    \right]\delta\vec{x} \ .
\end{align}
Because this is true for \emph{any} variation $\delta\vec{x}$ about a path $\vec{x}(t)$ that minimizes $S$, we find that the equation of motion is 
\begin{align}
    \frac{\partial L}{\partial \vec{x}}
    - 
    \frac{\D{}}{\D{t}}
    \frac{\partial L}{\partial \dot{\vec{x}}}  = 0 \ .
\end{align}
For the single particle Lagrangian above, this indeed gives
\begin{align}
    m\ddot{\vec{x}} = V'[\vec{x}] \ .
\end{align}


\section{Linear Algebra of Variations}

If you are like me, you remember the first time you learned the variational principle. It is so unusual compared to Newtonian mechanics that you may even remember the second and third time you learned the variational principle---because the first time didn't sink in. Why would it? Lagrangian mechanics is the first time you meet a new mathematical technique in physics: the \emph{calculus of variations}.

In case the usual treatment carries too much of an air of mystery, we can strip it down to its simplest componentnts by working in discretized time. In contrast to the calculus of variations, let us call this the \emph{linear algebra of variations}.\sidenote{There is no such phrase, nor is this actually different from the traditional treatment. We simply want to understand the calculus of variations from the perspective of linear algebra.} In this section, we first give a lightning review of the calculus of variations, then rephrase it in the more prosaic language of linear algebra.

\paragraph{Review of the calculus of variations}
In ordinary calculus, we vary some independent variable $x$ which may take any value in $\RR$. We search for extrema of the dependent variable, say $f(x)$, by looking at stationary points where $\delta f(x) = f(x+ \delta{x}) - f(x) = 0$. In the {calculus of variations}, we vary a separate independent variable $x(t)$ for \emph{each value of $t$}. That means that $x(0~\text{sec})$ and $x(1.2~\text{sec})$. There are a \emph{continuum} of independent variables, $x(t)$, ``indexed'' by the continuum paramter $t$. Of course, this is simply the path of the particle: for each time $t$, $x(t)$ tells you the position of the particle. The action functional $S[x(t)]$ is a number associated to each path.

In this calculus of variations, the analog of derivatives\sidenote{Math purists, forgive the abuse of the differential here. Just imagine $\D{x} \defeq \Delta x$ in the $\Delta x \to 0$ limit.} 
\begin{align}
    \D{f}(x) = f(x+\D{x}) - f(x) 
\end{align}
is a variation,
\begin{align}
    \delta S[x(t)] = 
     S[x(t)+\delta x(t)] - S[x(t)] 
     \ .
\end{align}
Ordinarly in calculus we write $\D{f}/\D{x}$, but for now I do not want to dig into what the `denominator' of $\delta S[x(t)]/\delta x(t)$ means: unlike $\D{x}$ which physicists like to think of as the infinitesimal limit of $\Delta x \to 0$, $\delta x(t)$ is a function. What does it mean to ``divide by a function?''\sidenote{The answer is that there is an implicit sum over indices. To see this, first examine the case where $f$ depends on two variables.} For what it is worth, the idea that $\D{f}/\D{x}$ is a fraction is already something that we do that causes mathematicians anxiety.\sidenote{This is formally true in an extension of standard analysis called \emph{nonstandard analysis} wherein infinitsimal limits are replaced by \emph{hyperreal numbers} that encode the meaning of `infinitseimal.' I do not know of any example where thinking of $\D{f}/\D{x}$ as a fraction leads to a logical inconsistency---but you can find some insightful discussion online.}

There is also an analog of integrals. In ordinary calculus, the integral is a sum over infinitesimal changes with respect to successive changes in the independent variable. For example, if $h(x) = f'(x)$,
\begin{align}
    \int_{x_0}^{x_1} \D{x} \, h(x)
    =
    \int_{x_0}^{x_1} \D{x} \, \frac{ \D{f} }{ \D{x} } 
    = f(x_1) - f(x_0)
    \ .
\end{align}
Observe that we seem to be able to `cancel out' the infinitesimals $\D{x}$. This is the fundamental theorem of calculus. In the calculus of variations, the analog of an integral is a \textbf{path integral} which is a sum over successive changes in the paths $\delta x(t)$:
\begin{align}
    \int \mathcal{D}x\, Z[x(t)] \ .
\end{align}
We denote the integral measure over all paths by $\mathcal{D}x$ to distinguish it from the single variable measure $\D{x}$---the path integral measure encodes an \emph{infinite} number of single variable measures.\sidenote{Observe that the integrand $Z[x(t)]$ is \emph{not} the action. The name $Z$ should make you think of the partition function.} Path integrals do not usually show up in classical mechanics, but they are a key part of the Lagrangian formulation of quantum mechanics.

With all that being said, let us convert the above machinery into something much more mundane: the linear algebra of variations. The key idea is this: rather than dealing with an \emph{infinite continuum} of independent variables $x(t)$, let us deal with a discrete set of variables that we may arrange as vectors $\ket{x}$. This is analogous to our ``histogram basis'' of function space. The vector space may still be infinite dimensional,\sidenote{Or you can always take a finite limit, say with periodic boundary conditions.} but the discrete spacing means that our differential operators are now naturally interpreted as large matrices rather.




\paragraph{Discrete Time} To avoid doing calculus, we simply replace the continuous time variable $t$ with a discrete lattice of specific times $t_i$ where $t_{i+1} = t_i + \Delta t$ for some finite\sidenote{I have a colleague who finds the word `finite' here rather odd. Usually `finite' is in contrast to infinite, but here we use `finite' in contrast to \emph{infinitesimal}. I refer to Jim Holt's essay on infinitesimals in \emph{When Einstein Walked with G\"odel} for a bit of the history of this idea.} and constant $\Delta t$.

This means that the trajectory of a particle through time, $x(t)$ is now reduced to a discrete list of numbers $x^i$ corresponding to $x(t_i)$. This list of numbers may be infinite, but they are countably infinite. Because we have discretized, there is nothing `in between' $t_i$ and $t_{i+1}$. Instead, we say that physics tells us how $x^{i}$ and $x^{i+1}$ are related to each other. In Newtonian mechanics, we would say that you simply have to know $x^i$ and some notion of the velocity, say $\Delta t\inv(x^i-x^{i-1})$, and you could then deduce $x^{i+1}$ using the discrete version of $F=m\ddot{x}$.
\begin{exercise}
If you have not done this before, it is a worthwhile exercise to write pseudocode for how a computer would determine the trajectory of an asteroid under the influence of a gravitational force over small discrete time steps. You can pretend that you are an engineer in \acro{NASA} in the 1960s. Even as someone who does not often write `honest' code, I highly recommend the classic \emph{Numerical Recipes 3rd Edition: The Art of Scientific Computing}~\autocite{10.5555/1403886} to go deeper into this topic.
\end{exercise}

The principle of least action is rather different from Newtonian mechanics. The action is a functional of the state vector, $S[x^i]$.\sidenote{In fact, you can think of $S[x^i] = \la \vec{x}, \vec{x}\ra$ for an appropriately defined inner product.} The equation of motion comes from minimizing $S$ for all possible values of $x^i$. Because $S$ is an integral over the Lagrangian, and because the Lagrangian's kinetic term relates the values of $x^i$ to that of $x^{i+1}$, we see that minimizing $S$ ends up relating the position of the particle at consecutive times.
\begin{align}
    S_\text{kin} = 
    \sum_{i} \Delta t \,
    \frac{1}{2} 
    \left(
        \frac{ x^{i+1} - x^i }{ \Delta t }
    \right)^2 \ .
\end{align}
The continuum limit $\sum_i \Delta t \to \int \D{t}$ should be clear. The variational principle is the statement that the classical path is the one for which $S$ is minimized over all possible positions at all possible (discrete) times. This variation is simply:
\begin{align}
    \delta S = \frac{\partial{S}}{\partial{x^i}}\delta x^i
    = 
    \frac{\partial{S}}{\partial{x^1}}\delta x^1
    +
    \frac{\partial{S}}{\partial{x^2}}\delta x^2
    + 
    \cdots \ .
    \label{eq:S:partial}
\end{align}
When $S$ is extremized, $\delta S = 0$ for \emph{any} set of variations $\delta x^i$.  Observe that we use the partial derivative notation here: $S$ is no longer strictly a \emph{functional}, but just a function of a large number of independent variables $x^i$. 
\begin{exercise}
Show that for a theory with no potential energy, the $\delta S_\text{kin}=0$ for arbitrary $\delta x^i$ implies that nearest [temporal] neighbor states are equal $x^i = x^{i+1}$. 
\end{exercise}

Let us be clear that in this section the particle moves in a single dimension. The index $i$ on $x^i$ refers to a time stamp, not a spatial direction.\sidenote{We'll get there, don't worry. For now, you can appreciate why it makes sense that in relativity, space and time are treated on similar footing.} The vector 
\begin{align}
    \ket{x} &= x^i \ket{i} 
\end{align}
is a description of all positions of $x(t_i)$ over the temoral history of the particle. The coefficients $x^i$ can take any conitnuous value, but the time slices $t_i$ are discrete.\footnote{The upper index on $x^i$ is really a vector index in that it is understood to contract with, say, basis kets. The lower index on $t_i$ is simply a discrete label.} The basis vectors $\ket{i}$ refer to a unit displacement of the particle in the positive direction at time $t_i$.
\flip{Insert graphic here. Make a discrete time diagram.}




\paragraph{Discretized Action Principle}
We would like to determine the function, $\ket{x}=x(t)$ that minimizes $S[x(t)]$. We may write $S$ as a sum over a Lagrangian $L$. 
\begin{align}
    S &= \sum_i \Delta t \, L[x^i, \dot x^i]  \ .
    % &
    % \ket{x} &= x^i \ket{i}
    % &
    % \ket{\dot{x}} &= \dot{x}^i \ket{i} \ ,
\end{align}
The Lagrangian itself is a function of $t$---that is, it depends on the discrete time index $i$.
Ah, but wait a moment: what do we make of the $\dot x^i$? The dot is a time derivative, and our goal here was to \emph{not do calculus}. We need to pass to the discretized form of the derivative---recall our `histogram basis' introduction to function spaces. We would like to write $\dot x$ as matrix representation of the time derivative acting on the $x$ vector:
\begin{align}
    \ket{\dot{x}} = D\ket{x} \ .
\end{align}
But $D$ itself is ambiguous: we need to decide whetehr $D$ is the \emph{forward} or the \emph{backward} derivative,
\begin{align}
    D_+ \ket{f} 
    &\equiv \sum_i \frac{ f^{i+1} - f^i }{\Delta t} \ket{i}
    \\
    D_- \ket{f} 
    &\equiv \sum_i \frac{ f^{i} - f^{i-1} }{\Delta t} \ket{i} \ .
\end{align}
For our purposes, let us choose $D=D_+$. We drop the explicit index for now when the difference is innoucous, but we restore it when we have to make use of the non-Hermiticity of the derivative, $D_\pm^\dag = - D_\mp$. The definitions above are subject to boundary conditions which we leave implicit.\sidenote{Assume periodic boundary conditions if ambiguity makes you anxious.} 
\begin{exercise}
Confirm that all of our discussion in the rest of this section holds equivalently if we had chosen $D=D_-$.
\end{exercise}

Armed with a definition for the discretized derivative, we may explicitly write out the Lagrangian for the kinetic term:
\begin{align}
    L[x^i]_\text{kin} &= \frac{1}{2} , (Dx)^i (Dx)^i 
    &
    \text{no sum over }i \ .
\end{align}
Eh, this is something of a cumbersome notation. The vector components $(Dx)^i$ are the coefficients of $D\ket{x} = (Dx)^i\ket{i}$, but here we have repeated upper indices\sidenote{This is already a sin!} and a reminder that there is no sum. It is much cleaner to write the action directly:
\begin{align}
    S_\text{kin} = \frac{m}{2} g_{ij} (Dx)^i (Dx)^j \ ,
    \label{eq:S:kin:discrete}
\end{align}
where now we have a completely scalar quantity with indices contracted. In fact, we see that $S_\text{kin} = (m/2)\la Dx, Dx \ra$. We take the metric to be Euclidean so that $g_{ij} = 1$ if $i=j$ and zero otherwise. For the sake of preserving the index structure, we keep writing it explicitly as $g_{ij}$.\sidenote{It is dangerous to write some ``Kronecker'' $\delta_{ij}$ because the object does not generically transform like a (0,2)-tensor.} 
% 
Observe that unlike $L_\text{kin}$, $S_\text{kin}$ already has a sum over all time slices. 


\paragraph{Equation of motion for the free particle, linear algebra version} A free particle is one which has no potential.\sidenote{There is something poetic here along the lines of \emph{The Big Lebowski}.} This means the action is $S=S_\text{kin}$. You already know that the equation of motion in the continuum case is $\ddot{x}(t)=0$. Let us see how we get this in the discretized case.  The key insight is that even though the time slices $t_i$ are discrete, the components $x^i$ are continuous. Thus we may ``do calculus'' on these objects as usual. In this way, the central objects are the \emph{partial derivatives} in \eqref{eq:S:partial}.

% There is a trick to make the steps conceptually clearer, even though it makes everything a bit clunky. Since standard textbooks do everything in a slick way, let's do the clunky derivation. 

Let us start by rewriting $S$ in the form $\frac{1}{2}\la x | A x \ra$, where $A$ is some self-adjoint\sidenote{$A^\dag = A$} matrix. In fact, $A$ is real so that $A^\trans = A$. Why would we want to do this? 
\begin{theorem}\label{thm:quadratic:action}
If $S=\frac{1}{2} \la x | A x \ra$, then the variation of $S$ is
\begin{align}
    \frac{\partial S}{\partial x^i}
    = 
    g_{ij}A\aij{j}{k}x^k \ .
    \label{eq:variation:of:quadratic}
\end{align}
\end{theorem}
\begin{proof}
The action is
\begin{align}
    S= \frac{1}{2}  g_{\ell j}x^\ell A\aij{j}{k}x^k \ .
\end{align}
The product rule from calculus gives
\begin{align}
    2 \frac{\partial S}{\partial x^i}
    &= 
    g_{\ell j}
    \frac{\partial x^\ell}{\partial x^i} A\aij{j}{k}x^k
    + 
    g_{\ell j} x^\ell A\aij{j}{k}
    \frac{\partial x^k}{\partial x^i}
    \\
    &=
    g_{\ell j}
    \delta^\ell_i A\aij{j}{k}x^k
    + 
    g_{\ell j} x^\ell A\aij{j}{k}
    \delta^k_i 
    \label{eq:discrete:2:partial:s:partial:xi:int}
    \ .
\end{align}
The first term on the right-hand side is 
$g_{i j} A\aij{j}{k}x^k$. This is half of the final answer \eqref{eq:variation:of:quadratic}. The second term had better be the same as the first term. To see this, tecall the definition of the adjoint for real matrices, \eqref{eq:adjoint:def}: $(A^\dag)\aij{i}{j} = g_{jk}A\aij{k}{\ell}g^{\ell i}$. Then the self-adjoitness of $A$ gives
\begin{align}
    A\aij{j}{k} = (A^\trans)\aij{j}{k}
    = g_{km}A\aij{m}{n}g^{nj} \ .
\end{align}
Plugging this into the second term of \eqref{eq:discrete:2:partial:s:partial:xi:int} gives
\begin{align}
    g_{\ell j} x^\ell A\aij{j}{k}
    \delta^k_i 
    &=
    g_{\ell j} x^\ell 
    g_{km}A\aij{m}{n}g^{nj}
    \delta^k_i 
    =
    g_{im}A\aij{m}{n}
    g^{nj}
    g_{\ell j} 
    x^\ell 
    =
    g_{im}A\aij{m}{\ell}
    x^\ell \ .
\end{align}
Up to the names of dummy indices, this indeed matches the first term of \eqref{eq:discrete:2:partial:s:partial:xi:int}. Changing these dummy indeices gives the result, \eqref{eq:variation:of:quadratic}.
\end{proof}

What remains, then, is to show how to write \eqref{eq:S:kin:discrete} in the form of Theorem \ref{thm:quadratic:action}. This, in turn, is straightforward from the definitino of the adjoint. Recall that everywhere we wrote $D$ actin on a vector we really meant the \emph{forward} derivative, $D_+$. Then thea ction is
\begin{align}
    S_\text{kin} = \frac{m}{2} \la D_+x , D_+x \ra  \ .
\end{align}
We also know that the adjoint of the forward derivative is the backward derivative, $D_+^\dag =- D_-^\dag$. By the definitin of adjoint, this means that 
\begin{align}
    \la D_+x , D_+x \ra = 
    - \la x , D_- D_+x \ra \ .
\end{align}
Nowe we recognize that $D^2 \equiv D_- D_+$ \emph{is} a self-adjoint operator. We may thus identify $A =- \frac{m}{2}D^2$ in Theorem \ref{thm:quadratic:action} \ . By the result of that theorem, we find that
\begin{align}
    \frac{\partial S_\text{kin}}{\partial x^i}
    = 
    - \frac{m}{2} g_{ij} (D^2)\aij{j}{k}x^k \ .
\end{align}
This is the equation of motion for the free particle. The continuum version is
\begin{align}
    \frac{\delta S_\text{kin}}{\delta x(t)}
    = - \frac{m}{2}\ddot x(t) \ .
\end{align}
\begin{exercise}
Walk through this same derivation in the continuum limit. It may help to write the action as:
\begin{align}
    S_\text{kin} = 
    \int \D{t}\,
    \frac{ m }{ 2 }
    \left[ 
        \frac{ \D{y} }{ \D{t} }
        \frac{ \D{z} }{ \D{t} }
    \right]_{y(t), z(t) =x(t)} \ .
\end{align}
We do this to help keep track of terms in the product rule:
\begin{align}
    \frac{ \delta S_\text{kin} }{ \delta x(t') }
    &=
    \int \D{t}\,
    \frac{ m }{ 2 }
    \left[ 
        \frac{\delta}{\delta y(t')}
        \left(\frac{ \D{y} }{ \D{t} }\right)
        \frac{ \D{z} }{ \D{t} }
        +
        \frac{ \D{y} }{ \D{t} }
        \frac{\delta}{\delta z(t')}
        \left(\frac{ \D{z} }{ \D{t} }\right)
    \right]_{y(t), z(t) =x(t)} \ .
\end{align}
Here the $\delta/\delta x(t')$ is a variation at some fixed $t'$ in the same way that $\partial/\partial x^i$ is a varation of the path at some fixed $t_i$. Now the challenge is: how do we vary a term like $\D{y}/\D{t}$ with respect to $\delta y(t')$? We can integrate by parts so that
\begin{align}
    \int \D{t}\,
    \frac{\delta}{\delta y(t')}
        \left(\frac{ \D{y} }{ \D{t} }\right)
        \frac{ \D{z} }{ \D{t} }
    =
    -
    \int \D{t}\,
    \frac{\delta}{\delta y(t')}
        y(t)
        \frac{ \D[2]{z} }{ \D{t}^2 } 
        \ .
\end{align}
Then use the fact that $\delta y(t)/\delta y(t') = \delta(t-t')$ in the same way that $\partial x^i/\partial x^j = \delta^i_j$:
\begin{align}
    \int \D{t}\,
    \left[\frac{\delta}{\delta y(t')}
                \left(\frac{ \D{y} }{ \D{t} }\right)
                \frac{ \D{z} }{ \D{t} }\right]_{y,z = x}
    =
    - \frac{ \D[2]{x} }{ \D{t'}^2 } \ .
\end{align}
Fill in the remaining steps as needed.
\end{exercise}
The `linear algebra' approach to solving the quadratic is \emph{not} the one you see in analytical mechanics, but it is more convenient in field theory.





\paragraph{Equation of motion for the free particle, calculus version}
Let us redo the free particle with discretized time, but this time let us do it the way that an analytical mechanics textbook might do it. We start from \eqref{eq:S:kin:discrete} and rewrite it in terms of an intermediate variable $y^i$,
\begin{align}
    S_\text{kin} &= \frac{m}{2}g_{ij} y^i y^j
    &
    y^i\defeq D\aij{i}{k} x^k \ .
\end{align}
This variable $y(t)$ is what we typically call the velocity $\dot x(t)$.\sidenote{Remember we start by pretending that the velocity is independent of the path, then we remember that it is not.} By the chain rule, we have
\begin{align}
    \frac{ \partial S_\text{kin} }{ \partial x^i }
    =
    \frac{ \partial S_\text{kin} }{ \partial y^j }
    \frac{ \partial y^j }{ \partial x^i } \ .
\end{align}
The factors are simple:\sidenote{Fill in the steps.}
\begin{align}
    \frac{ \partial S_\text{kin} }{ \partial y^j }
    &= m g_{jk} y^k
    &
    \frac{ \partial y^j }{ \partial x^i }
    &=
    D\aij{j}{i}
    \ .
\end{align}
The second equation makes sense in the discrete case where $D$ is just some matrix, but it does not seem to make any sence in the continuum case where $D\to \D{}/\D{t}$ is a differential operator---what is it acting on? In order to make our expression sensible in the continuum limit, we should make sure $D$ is acting on something. Fortunately, index notation means we can move factors around:
\begin{align}
    \frac{ \partial S_\text{kin} }{ \partial x^i }
    &=
    m g_{jk} y^k
    D\aij{j}{i}
    =
    % m
    % D\aij{j}{i} g_{jk} y^k
    % =
    m 
    \delta^\ell_i
    D\aij{j}{\ell} g_{jk} y^k
    = 
    m g_{im}\, g^{m\ell} D\aij{j}{\ell} g_{jk} \, y^k \ .
\end{align}
where we use $\delta^\ell_i = g_{im} g^{m\ell}$. Now we recognize the transpose: 
\begin{align}
    g^{m\ell} D\aij{j}{\ell} g_{jk} \equiv (D^\trans)\aij{m}{k}
    \ .
\end{align}
Because we have assumed $D=D_+$, we  have $D^\trans = - D_-$. As a result, we have
\begin{align}
    \frac{ \partial S_\text{kin} }{ \partial x^i }
    = - m g_{im} (D_-)\aij{m}{k} (D_+)\aij{k}{j} x^j \ .
\end{align}
As before, we know that the product $D_-D_+ = D^2$ is the self-adjoint Laplacian operator so that we recover the usual $-m\ddot{x} = 0$ equation of motion. 

\begin{exercise}
Show that the continuum limit of this procedure is exactly what we do in Section~\ref{sec:lagrange:spring:theory}
\end{exercise}



\section{Discrete Time Harmonic Oscillator}

The harmonic oscillator action has a potential term that is quadratic in the position variable. In the discrete time limit,
\begin{align}
    S = 
        \frac{m}{2} g_{ij}(Dx)^i(Dx)^j 
        - \frac{k}{2} g_{ij} x^i x^j \ .
\end{align}
The relative minus sign ensures that there is a stable minimum energy configuration.
\begin{exercise}
Using complete sentences, explain why this is the most generic form of an action that is quadratic in $x$. Use homogeneity and isotropy in time to argue why there can be no tensor structures other htan $g_{ij}$.
\end{exercise}

Appreciate that the \emph{quadratic} part of the action gives a linear equation of motion. This is obvious because varying a quadratic action with respect to $x^i$ gives terms that are linear in the position variable. This means that higher-order terms, e.g.\ trilinear terms like $(x^i)^3$, must be treated separately.\footnote{Often we treat these as a small correction to the quadratic action. This amounts to saying that everything is a perturbation on a harmonic oscillator. The perturbations can be treated systematically using a technique called Feynman diagrams.}
% \begin{exercise}
% Confirm that \eqref{eq:HO:action:discretized} is the discretization of the action for a simple harmonic oscillator. For example, show that in the continuum limit you recover the simple harmonic oscillator action. 
% \end{exercise}
% 
% 

% \flip{This next exercise is probably wrong: I think the point is that $D_+^\dag = D_-$ so that we can define the forward derivative for vectors and the backward derivative for forms. The action takes the form $(D_+q)^2$ so that the equation of motion is $D_-D_+ q=0$.}
% \begin{exercise} \label{ex:integrate:kinetic:term:by:parts}
% If $x(t)$ is the position of a particle of mass $m$, then its kinetic term is $m\dot x(t)/2$. One way to derive the equation of motion for a theory is to integrate this term by parts, see Section~\ref{eq:variational}. Let us see how this works. Rather than $\dot x^2$, let us consider $(D_+f)(D_-g)$. We write $f$ and $g$ to make it easier to keep track of terms, but for the specific case of the kinetic term we know $f=g=x(t)$. 

% The integral over time corresponds to a sum over the time slice (histogram) index $i$ of $f$ and $g$:
% \begin{align}
%     \int dt\, \dot x(t)^2 \to \sum_i \Delta x \, (D_+f)^i(D_-g)^i \ .
% \end{align}
% Show that this sum takes the form
% \begin{align}
%     \sum_i \Delta x \, (D_+f)^i(D_-g)^i 
%     = -\sum_i f^i (D^2 g)^i 
%     + f^{N+1}(g^N-g^{N-1})
%     - f^1(g^1 - g^0) \ .
% \end{align}
% I was sloppy with the limits of the sum above. What values does $i$ range over in that sum?
% % HINT:
% \acro{Partial answer:}
% Here's the pattern:
% \begin{align}
%     D_+fD_-g &= 
%     f^2g^1 - f^2g^0 - f^1 g^1 + f^1 g^0\\
%     &+ f^3 g^2 - f^3 g^1 - f^2 g^2 + f^2 g^1 \\
%     &+ f^4 g^3 - f^4 g^2 - f^3 g^3 + f^3 g^2 \\
%     &+ \cdots \\
%     & + f^{N+1} g^N - f^{N+1} g^{N-1} - f^N g^N + f^N g^{N-1}  \ .
% \end{align}
% \end{exercise}
% 
% 

% Recalling Exercise~\ref{ex:integrate:kinetic:term:by:parts}, we may integrate the derivative term by parts to obtain

From the previous section, we are now confident that we can integrate the first term by parts---in fact, we saw how to interpret this in the discrete limit two ways---so that
\begin{align}
    S = 
    -
    \frac{m}{2} g_{ij}x^i (D^2x)^j 
    - \frac{k}{2} g_{ij} x^ix^i
    \equiv
    -\frac{1}{2}x^i Q_{ij} x^j
    \ ,
    \label{eq:HO:action:discretized:integrated}
\end{align}
where we have defined the `quadratic term' operator
\begin{align}
    Q_{ij} = 
    m\, g_{ik}(D^2)\aij{k}{j} 
    + 
    k\, g_{ij} 
    = Q_{ji} \ .
    \label{eq:Qij}
\end{align}
\begin{exercise}
Confirm that $Q_{ij}$ is symmetric, $Q_{ij} = Q_{ji}$. 
\end{exercise}
\begin{exercise}
What is the \emph{support} of \eqref{eq:Qij}? That is, what is the condition on $i$ and $j$ to find a non-zero $Q_{ij}$?
\end{exercise}

The principle of least action posits that the classical trajectory $x(t)$---encoded in the discretized $x^i$ coefficients---is the configuration that extremizes $S$. To do this, we vary each $x^i$ independently. This is what the functional variation $\delta x(t)$ means in the continuous limit. The extremization corresponds to the $N$ equations
\begin{align}
    \frac{\partial{S}}{\partial{x^1}} &= 0
    &
    \frac{\partial{S}}{\partial{x^2}} &= 0
    &
    \cdots&
    &
    \frac{\partial{S}}{\partial{x^N}} &= 0 \ .
\end{align}
where relevant, these are subject to the boundary conditions. 
% 
Of course, the $x^i$ are not independent: we know that the kinetic term relates the displacement at time $i$ to the neighboring displacements at time $i\pm 1$.
% 
Rather than solving each of these one by one, we can solve the equation for a generic $x^i$:
\begin{align}
    -\frac{ \partial{S} }{ \partial{x^i} } &= 
    \frac{ \partial }{ \partial{x^i} }
    \left[
    \frac{1}{2} Q_{ii}(x^i)^2 + \sum_{j\neq i} \frac{1}{2}\left(Q_{ij} x^ix^j + Q_{ji}x^jx^i + \right) 
    \right] \\
    &= 
    Q_{ii}x^i+ \sum_{j\neq i} Q_{ij}x^j \ .
\end{align}
where there is no sum over repeated $i$ indices and we have explicitly written the sum over $j\neq i$. In the last term we have used $Q_{ij} = Q_{ji}$. Expanding this expression gives
\begin{align}
    0&=
    -\frac{\partial{S}}{\partial{x^i}}
    \\
    &=
    -2\frac{m}{\Delta t^2} x^i + k x^i 
    + \frac{m}{\Delta t^2}x^{i+1} 
    + \frac{m}{\Delta t^2}x^{i-1}
    \\
    &= 
    g_{ij}\left[m(D^2x)^j + k x^j\right] \ .
\end{align}
This gives us the discretized equation of motion:
\begin{align}
    (D^2x)^j = - k x^j \ ,
\end{align}
which we recognize is the discretized Hooke's law $m\ddot{x} = -k x$. 

What is significant about this approach is that we never invoked the standard Euler--Lagrange trick of first treating $x$ and $\dot{x}$ as independent variables and then `remembering' that they are related by a time derivative when it was convenient. We took a more direct path where we independently varied each piece of the path in time, $x^i$, while treating a differential operator $Q_{ij}$ as a matrix. We noticed that the $x^i$ are \emph{coupled} by the kinetic term, which tells us that wiggles in $x^i$ will affect the variation of $x^{i+\pm 1}$. In quantum mechanics, this approach is known as the path integral formulation.

\begin{example}
Having established the variational principle approach to the equation of motion ,we may repeat the derivation of Hooke's law by starting with the continuum limit. The action is
\begin{align}
    S&= \int_{t_0}^{t_1} dt\, \frac{m}{2}\dot x(t)^2 - \frac{k}{2}x(t)^2
    \\&
    = -\int_{t_0}^{t_1} dt\, \frac{1}{2}x(t)
    \left[m\left(\frac{d}{dt}\right)^2 + k\right]x(t)
    \\&    
    \equiv -\int_{t_0}^{t_1} dt\, \frac{1}{2} x(t) \mathcal O x(t) \ .
\end{align}
The classical equation of motion is thus
\begin{align}
    \mathcal Ox(t) = \left[m\left(\frac{d}{dt}\right)^2 + k\right]x(t) = 0 \ m \ ,
\end{align}
or in other words: $\ddot x(t) = -kx(t)$.
\end{example}
We have implicitly assumed that the boundary terms vanish.

On physical principles, the kinetic term is typically proportional to a term with two time derivatives. You can argue this based on spacetime symmetry. This means that one must perform an integration by parts on this term to write the equation of motion operator $\mathcal O$. This may give some insight on the meaning of the relative minus sign between the kinetic and potential energy terms in the Lagrangian. 




\section{Coupled Harmonic Oscillators}


Let's take stock of what we have accomplished. We considered a particle that can move in one direction of space subject to a harmonic oscillator potential. We discretized time to dissect the meaning of the variational principle. There is a trivial extension if the particle can move in multiple spatial directions: simply replace $x(t) \to \vec{x}(t)$ where $\vec{x}$ is a \emph{spatial} vector.\sidenote{Do not confuse this with $\ket{x}$ which is a vector with respect to the lattice of time slices.} 
\begin{exercise}
Work out the equation of motion for a particle in three dimensions of space subject to a harmonic oscillator potential. The spring constant $k$ is upgraded to a symmetric matrix $k\aij{i}{j}$, which you may take to be $k\delta^i_j$ for simplicity. Be sure to distinguish between $\partial/\partial x^i$, $\partial/\partial y^i$, and $\partial/\partial z^i$.
\end{exercise}
\begin{exercise}
What happens when $k\aij{i}{j}$ is not proportional to $\delta^i_j$?
\end{exercise}

We now want to move from single particles to fields. 
A particle has a single position at each time. Actually, we need to be more precise with our language. Rather than talking about `positions' of a particle, we should talk about a \emph{displacement from equilibrium.} So rather than talking about $x(t)$, we write $q(t)$.\sidenote{For example, we may not be tracking an actual position of a particle, but the temperature at a single point.} In contrast, a \textbf{field}\index{field} is an object that has some displacement at each time \emph{and at each position}. 

In the discretized limit, rather than just having a lattice of time slices, we have a lattice of space and time. For a given $x_a$ and $t_i$. The basis vectors are $\ket{a,i} = \ket{a}\otimes \ket{i}$: a unit displacement at position $x_a$ and time slice $t_i$. A state is $\ket{q} = q^{ai}\ket{a,i}$. 
% 
Rather than going through everything with both discrete space and discrete time, let's be lazy and return to continuum time. 

Consider an array of $N$ harmonic oscillators. The displacement at time $t$ of each oscillator $i$ with respect some baseline value to be $q_i(t)$. If each harmonic oscillator is independent, then each one would have a separate equation of motion---this is essentially $N$ copies of the single harmonic oscillator. Let us instead consider the case where
\begin{enumerate}
    \item The point $q_i(t)=0$ is not special and does not represent an equilibrium value. In fact, let us throw out the $V[q_i]=\frac{1}{2}kq_i^2$ potential energy term and replace it with something else.
    \item Let us couple the harmonic oscillator to the positions of its \emph{nearest neighbors}. The nearest neighbors of the $q_i$ oscillator are the $q_{i\pm 1}$ oscillators. 
\end{enumerate}
The potential energy of the system looks like
\begin{align}
    V[q_1,\cdots q_N] &=
    \cdots
    \frac{1}{2}k(q_{i+1}-q_i)^2 +
    \frac{1}{2}k(q_{i}-q_{i-1})^2 +
    \cdots \ .
\end{align}
This simply means that the locations of the $i^\text{th}$ oscillator's neighbors will pull the $i^{\text{th}}$ oscillator in one direction or the other. There is \emph{no} potential that pulls $q_i$ to its origin: $(k/2)q_i^2$. 


The Lagrangian  for this system is
\begin{align}
    L[q_1,\cdots, q_N]
    &= 
    \sum_i \frac{1}{2} m \dot q_i^2 - \frac{1}{2}k (q_{i+1}-q_i)^2 \ .
\end{align}
You can work out the equation of motion simply by varying the action, $S=\int dt\,L$. For example, the equation of motion for the $q_i(t)$ oscillator is
\begin{align}
    m\ddot{q}_i(t) = -k(q_{i}-q_{i-1}) -k(q_{i}-q_{i+1}) \ .
\end{align}
We observe that the spring constant terms push to decrease the difference between $q_i$ and its nearest neighbors. The force from the neighboring springs are in the direction to reduce the difference. In the continuum limit,
\begin{align}
    m\ddot{q}(t,x) = k\frac{\partial^2{q(t,x)}}{\partial{x}^2} \ .
\end{align}
We can write this as the d'Alembertian,
\begin{align}
    0 = m \partial^2 q \equiv m\left(\partial_t^2 - c^2 \partial_x^2\right) q(t,x) \ ,
\end{align}
where we recognize the speed of propagation $c^2 = k/m$.





\section{Feynman}
% Classical to quantum
% \flip{Classical to quantum, not just minimum, but the entire ensemble. }

The \textbf{path integral}\index{path integral} is something that shows up in the `sum over histories' approach to quantum mechanics. It is used extensively in quantum field theory because it gives a powerful way to impose constraints on a quantum system using Lagrange multipliers. What is not always obvious is that the path integral can be a more intuitive interpretation of the standard Euler--Lagrange approach. 

\paragraph{Sum Over Histories}
% 
\flip{To be filled in. Please see the example in class of the $N$-slit experiment and the proposal that each path contributes $\e^{iS}$.}

\begin{newrule}[Sum over histories]
Given some observed initial state $\ket{\text{in}}$ and some observed final state $\ket{\text{out}}$, the amplitude to go from $\ket{\text{in}}$ to $\ket{\text{out}}$ is
\begin{align}
    \mathcal M(\text{in}\to\text{out})= 
    \sum_\textnormal{path} \e^{iS_\textnormal{path}} \ ,
\end{align}
where the sum is over all `paths' (histories) that permit the initial state to evolve into the final state. $S_\textnormal{path}$ is the action for a given path---the same action that you are used to from classical mechanics, the time integral over the Lagrangian, $S= \int \D{t}\, L$.
\end{newrule}

\paragraph{Path Integral}
Suppose a quantum state has initial and final positions\sidenote{We deliberate use the notation $q$ for position instead of $x$. This is to aid our transition to quantum \emph{field theory}. Old textbooks sometimes refer to \emph{second quantization}; it is precisely to distinguish the role of $q$, a field displacement, and $x$ a continuous field index indicating spacetime position.}
\begin{align}
    \ket{\textnormal{in}} &= \ket{q_0}
    &
    \ket{\textnormal{out}} &= \ket{q_T} \ .
\end{align}
These are labeled so that we observe $\ket{q_0}$ at time $t=0$ and we observe $\ket{q_T}$ at time $t=T$. The amplitude to go from $\ket{q_0} \to \ket{q_T}$ over the time interval $T$ is
\begin{align}
    \bra{q_T} \hat U(T) \ket{q_0} &= \bra{q_T} \e^{-iT\hat H} \ket{q_0} \ .
\end{align}
Here $\hat U(t)$ is the operator that evolves states over time $t$. As you know from quantum mechanics, the Hamiltonian $\hat H$, is the operator that enacts infinitesimal evolution in time. This means that we enact finite time translations by exponentiating the Hamiltonian.\sidenote{Please make sure you are familiar with this idea: Hermitian operators \emph{generate} infinitesimal changes to a system. The exponential of these Hermitian operators are unitary operators that enact finite transformations. The idea shows up over and over again in this class.}

There are two tricks to deploy at this step:
\begin{enumerate}
    \item Break the finite time evolution operator into a product of $(N+1)$ time evolutions:
    \begin{align}
        \hat U(T) = \hat U(\Delta t)\hat U(\Delta t)\cdots \hat U(\Delta t) \ .
    \end{align}
    It is \emph{obvious} that $\hat U(t)\hat U(s) = \hat U(t+s)$.
    \item Insert the identity as a sum over a complete set of intermediate states:
    \begin{align}
        \one &= \int \D{q_i}\, \ket{q_i}\!\bra{q_i} \ .
    \end{align}
    We insert a copy of the identity in-between each of the $U(\Delta t)$s. Thus the label $i$ runs from $1$ to $N$.
\end{enumerate}
\begin{exercise}
Confirm that $\hat U(t)\hat U(s) = \hat U(t+s)$. \textsc{Hint}: this is trivial, but convince yourself that it is trivial.
\end{exercise}

Applying these two tricks to our amplitude\sidenote{$\D[N]{\vec{q}} \defeq \D{q_1}\cdots\D{q_N}$} 
\begin{align}
    \bra{q_T} \hat U(T) \ket{q_0} 
    &= \int \D[N]{\vec{q}} \,
    \bra{q_T} \hat U(\Delta T) \ket{q_N}
    \bra{q_N}\hat U(\Delta T) \ket{q_{N-1}}\cdots 
    \\
    &=\int \D[N]{\vec{q}} \,
    \prod_i^{N+1} K_i \ , \label{eq:path:integral:time:evolution:product:of:Ks}
\end{align}
where the factor $K_i$ is a complex number:
\begin{align}
    K_i &\defeq \bra{q_{i}} \hat U(\Delta T) \ket{q_{i-1}} \ ,
\end{align}
where we define $\bra{q_{N+1}} = \bra{q_T}$. We have reduced the amplitude to go from our initial state to our final state to a product of $K_i$. The form of $K_i$ is simply that of the amplitude to go from $\ket{q_{i-1}}$ to $\ket{q_{i}}$ over a small time $\Delta t$. The integral over intermediate states $\D[N]{\vec{q}}$ tells us that we are summing over all possible configurations of those unobserved intermediate states. 

In order to evaluate $K_i$, let us remind ourselves of the form of the Hamiltonian in $U(\Delta T) = \exp(i\Delta T\hat H)$,
\begin{align}
    \hat H = \frac{\hat P^2}{2m} + V(\hat Q) \ ,
\end{align}
where we use the non-relativistic kinetic term $\frac{1}{2}m\dot q^2 = p^2/2m$ and an arbitrary potential $V(q)$ that only depends on the $q$ displacements, but not the momenta $p$. Because $\hat H$ is an operator, we write $\hat P$ for the momentum operator and $\hat Q$ for the position\sidenote{This `position' is not necessarily spatial position. In general it is a displacement of a field value.} operator. All Hamiltonians in ordinary field theory have this type of separation: there is a \emph{kinetic term} that depends on momentum and a \emph{potential term} that depends on displacements. When we pass to relativistic theories, the kinetic term converts into a form that includes spatial derivatives in a way that it is Lorentz invariant. 

Once again there is another silly trick. 
\begin{enumerate}
    \setcounter{enumi}{2}
    \item Insert the momentum-space representation of the identity,
    \begin{align}
        \one &= \int \Dbar{p}\, \ket{p}\!\bra{p} 
        &
        \Dbar{p}\defeq \frac{\D{p}}{2\pi} \ .
    \end{align}
    Observe that this has a relative factor of $2\pi$ compared to the resolution of the `position space' identity. This is the famous---and seemingly inscrutable---factor that shows up in the Fourier transform, see Appendix~\ref{app:Fourier}.\sidenote{If you spend time thinking about this, you may wonder why we do not split the $2\pi$ evenly between positions and momenta. This would certainly be symmetric, but this choice ends up being fairly convenient in physics. It means, for example, that the conjugate variable for time is essentially the angular momentum, $E=\omega$.}\footnote{See e.g.\ \url{https://physics.stackexchange.com/a/141737/166736}} 
\end{enumerate}
Because $\hat H$ contains a term that evaluates easily in momentum space and a term that evaluates easily in position space, we insert the complete set of momentum states into $K_i$ to allow the $\hat P$-dependent terms to hit a $\bra{p}$ while the $Q$-dependent terms hit a $\ket{q_i}$:
\begin{align}
    K_i &=
    \int \Dbar{p} \la {q_{i}} | p \ra \bra{p} \e^{-i\Delta t\frac{ \hat P^2}{2m}} \e^{-i\Delta tV(\hat Q)} \ket{q_{i-1}}
    \label{eq:path:integral:Ki:operator}
    \\ &=
    \int \Dbar{p} \, \e^{-i\Delta t\frac{ p^2}{2m}} \e^{-i\Delta tV(q_{i-1})}\,  \la {q_{i}} | p \ra \la p | q_{i-1} \ra  \ .
    \label{eq:path:integral:Ki:eigen}
\end{align}
The factors on the right are simply the Fourier components between the position and momentum eigenstates:
\begin{align}
    \la {q_{i}} | p \ra &= \e^{-ipq_i}
    &
    \la p | q_{i-1} \ra &= \e^{+ipq_{i-1}}
    \ .
\end{align}
Plugging this into $K_i$ gives
\begin{align}
    K_i &=
    \e^{-i\Delta tV(q_{i-1})} \int \Dbar{p} \, \e^{-i\Delta t\frac{ p^2}{2m} - i p(q_i - q_{i-1}) } \ .
\end{align}
If you do not worry\sidenote{%
    Will not worry. On the other hand, you may want to worry about whether the assumptions we made to solve these Gaussian integrals is still valid when there are complex coefficients. This scratches the surface of why some mathematicians are skeptical about whether anything we do in quantum field theory is well posed.} 
about the factors of $i$, then this integral is precisely in the form of a Gaussian with a source term. This has a known solution:
\begin{align}
    G_{a,J} \equiv \int \D{x}\; e^{-\frac{a}{2}x^2 + Jx} = \sqrt{\frac{2\pi}{a}}  e^{\frac{J^2}{2a}} \ .
    % \tag{\ref{eq:Gaussian:with:source}}
\end{align}
We identify
\begin{align}
    a &= \frac{i\Delta t}{m}
    &
    J&= -i(q_i - q_{i-1}) \ .
\end{align}
Plugging in this result gives
\begin{align}
    K_i &=
    \e^{-i\Delta tV(q_{i-1})} 
    \frac{1}{2\pi}
    \sqrt{\frac{2\pi m}{i \Delta t}}
    \e^{\frac{-(q_i-q_{i-1})^2}{2i\Delta t/m}}
    \\
    &=
    \sqrt{\frac{m}{i 2\pi \Delta t}}
    \e^{\frac{i}{2}m\frac{(q_i-q_{i-1})^2}{\Delta t^2}\Delta t}
    \e^{-iV(q_{i-1})\Delta t} 
    \\
    &= \mathcal N \e^{i\Delta t\, L[q_i]}
    \label{eq:path:integral:infinitesmimal:kernel}
     \ .
\end{align}
In the last line we have identified $(q_i-q_{i-1})/\Delta t = \dot q_i$ and observed that the argument of the combined exponential is simply $i\Delta t$ times the \emph{Lagrangian} at point $i$. The ugly square root is some complex constant factor that \emph{does not matter}. We simply write it as a normalization $\mathcal N$. 

Now we insert this into the amplitude to go from our initial state to our final state, \eqref{eq:path:integral:time:evolution:product:of:Ks}:
\begin{align}
    \bra{q_T} \hat U(T) \ket{q_0} 
    &= \mathcal N^{N+1} \int \D[N]{\vec{q}} \,
    \e^{i \sum_i \Delta t\, L[q_i]} \ .
\end{align}
Recalling that $i$ encodes `time slices,' we recognize that the argument of the exponential is simply an \emph{integral} of the Lagrangian,
\begin{align}
    i \sum_i \Delta t\, L[q_i] \to i\int \D{t}\, L[q(t)] \equiv iS[q(t)] \ ,
\end{align}
where we have recovered the action. In the continuum limit, it is conventional to replace the notation $\D[N]{\vec{q}}$ with $\mathcal{D}q(t)$ so that our final expression for the $\Delta t \to 0$ limit is
\begin{align}
    \bra{q_T} \hat U(T) \ket{q_0} 
    &= \mathcal N' \int \mathcal{D}q(t)\;
    \e^{i S[q(t)]} \ .
    \label{eq:path:integral:formalism}
\end{align}
Here $\mathcal N'$ is a re-packaging of the powers of $\mathcal N$. The integration over all trajectories $\mathcal {D}q(t)$ is called a \textbf{path integral}\index{path integral}.

\begin{exercise}
Why was it necessary for us to split $U(T)$ into several small steps? \textsc{Hint}: did we make any assumptions to go from \eqref{eq:path:integral:Ki:operator} to \eqref{eq:path:integral:Ki:eigen}? 
\end{exercise}

\paragraph{Recovering Lagrange}
% \flip{To be filled in. Discussion of phasor diagrams, steepest descent, and why the classical path dominates. Motivates the principle of least action as the $\hbar \to 0$ limit.}
The path integral formulation limits to the classical equation of motion in the limit $\hbar \to 0$. 
\begin{exercise}
Show that the $\hbar \to 0$ limit reduces to the standard Euler--Lagrange equation of motion for a classical system. You can argue heristically using a tool called \textbf{phasor diagrams}, though a more technical derivation will take the leading order terms in an expansion in $\hbar$.
\end{exercise}


\paragraph{Recovering Schr\"odinger}
This section is outside the primary scope of the class. However, I encourage you to at least skim it to see an alternative explanation of where the Schr\"odinger equations comes from. It should convince you that the sum over histories (path integral) construction gives you the same dynamics as before. 
% 
% 
We show that the Schr\"odinger equation appears as the single-particle, non-relativistic limit over the path integral formalism.\sidenote{This section draws from the excellent pedagogical article by David Derbes in the \emph{American Journal of Physics} (one of my favorite publications) in 1996. This artcile was selected as one of the Fermat's Library journal club articles.\footnotemark}\footnotetext{\cite{10.1119/1.18114}; \url{https://fermatslibrary.com/s/feynmans-derivation-of-the-schrodinger-equation}} The power of the path integral formalism is that it generalizes beyond non-relativistic quantum mechanics. Like with all new theoretical formalisms, it behooves one to confirm that the new formalism is mathematically equivalent to the established formalism in the appropriate limit.

Suppose we interpret $q(t)$ as the position of a non-relativistic particle at time $t$. Then in this basis, we can think of the amplitude $\la q_t | \hat U(t) | q_0 \ra$ as the wavefunction $\psi(q, t)$ for the particle to be in position $q$ at time $t$ given some initial state $\ket{q_0}$ at $t=0$. We can rewrite \eqref{eq:path:integral:time:evolution:product:of:Ks} as follows:
\begin{align}
    \psi(q,t) &= \la q_t | \hat U(t) | q_0 \ra 
    \\&= \int \D{q_N} \, K_{\!N} \int \D[N-1]{\vec{q}} \,
    \prod_i^{N} K_i 
    \\
    &= \int \D{q'} \, K_{\!N}\, \psi(q',t-\Delta t) \ .
\end{align}
We have relabeled $q_N \to q'$. It is the last displacement variable that we integrate over. In fact, for convenience,\sidenote{You do not have to do this step.} let us shift all of our time coordinates by $\Delta t$: 
\begin{align}
    \psi(q,t+\Delta t) 
    &= \int \D{q'} \, K_{\!N}\, \psi(q',t) \ .
\end{align}
This equation says that to go from the $(N-1)^\textnormal{th}$ step to the $N^\textnormal{th}$ step, you integrate the wavefunction at last position over all possible intermediate steps, weighted by the \emph{kernel},\sidenote{Kernels are just fancy names for distributions, which are in turn fancy names for functions that are integrated over.} $K_N$. We can now use the form of the kernel, \eqref{eq:path:integral:infinitesmimal:kernel}:
\begin{align}
    K_{\!N} &= \mathcal N \e^{i\Delta t\, L[q_i]}
    &
    \mathcal N = \sqrt{\frac{m}{2\pi i \Delta t}} \ .
    \label{eq:path:integral:normalization}
\end{align}
Let us now be careful with the scaling with the small time slice $\Delta t$:
\begin{align}
    i\Delta t L[q_i] = \frac{im}{2}\frac{(q_i-q_{i-1})^2}{\Delta t} - i\Delta t V(\bar q) \ .
\end{align}
Note that the kinetic term goes like $\Delta t^{-1}$, which is \emph{large}. The potential term scales linearly with $\Delta t$, so it is small. We write $V(\bar q)$ to mean the average value of $V$ in the interval $(q_{i-1}, q_i)$. For our purposes, $\bar q = q$ assuming a sufficiently smooth potential. This means we can Taylor expand the potential term, but not the kinetic term:
\begin{align}
    \psi(q,t+\Delta t) &= \mathcal N\int \D{q'}\; 
    \e^{\frac{im}{2}\frac{(q-q')^2}{\Delta t}} 
    \left[1 - i\Delta t \, V(\bar q)\right]\,
    \psi(q',t) \ .
\end{align}
We drop higher order terms in $\Delta t$.
Now we assume that $q= q' + \xi$ and make the plausible case that $\xi$ is small.\sidenote{To be clear, we integrate over even large values of $\xi$. The sin that I am sweeping under the rug is the idea that the rapidly varying oscillation in the $\e^{\textnormal{kinetic}}$ term will cause integrants with large values of $\xi$ to cancel. This is the steepest descent method in the method of phasors.} This means we can also replace $\bar q \to q$ to leading order.
We can then expand $\psi(q',t)$, where we shall suppress the $t$ argument so that spatial derivatives are clear:
\begin{align}
    \psi(q') = \psi(q) -\xi \psi'(q) + \frac{1}{2}\xi^2 \psi''(q) + \cdots \ .
\end{align}
To be clear: we now write $\psi'(q,t)$ to mean $\frac{\D{}}{\D{q}}\psi(q,t)$.
Plugging this in:
\begin{wide}
\begin{align}
    \psi(q,t+\Delta t) &= 
    \mathcal N\int \D{\xi}\; 
    \e^{\frac{im}{2}\frac{\xi^2}{\Delta t}} 
    \left[1 - i\Delta t \, V(q)\right]\,
    % \\
    % &\phantom{=
    % \mathcal N\int \D{\xi}\; 
    % }\times 
    \left[\psi(q,t) -\xi \psi'(q,t) + \frac{1}{2}\xi^2 \psi''(q,t)\right] 
    \\
    % &= 
    % \mathcal N\int \D{\xi}\; 
    % \e^{\frac{im}{2}\frac{\xi^2}{\Delta t}} 
    % \left[1 - i\Delta t \, V(\bar q)\right] \psi(q,t)
    % \\
    % &\phantom{=}+
    % \mathcal N\int \D{\xi}\; 
    % \e^{\frac{im}{2}\frac{\xi^2}{\Delta t}} 
    % \frac{1}{2}\xi^2 \psi''(q,t) \ .
    &= 
    \mathcal N\int \D{\xi}\; 
    \e^{\frac{im}{2}\frac{\xi^2}{\Delta t}} 
    \left[\psi(q,t) - i\Delta t \, V(q)\,\psi(q,t)
    + \frac{1}{2}\xi^2 \psi''(q,t)
    \right]
    \ .
 % \ . \nonumber
 \label{eq:psi:path:integral:schrodinger:intermediate}
\end{align}
\end{wide}
In the last line we removed the term linear in $\xi$ because it cancels against the rest of the integral, which is even in $\xi$. We also remove terms of order $\mathcal O(\Delta t\,\xi^2)$ since these are small. Observe that the $\xi$-independent terms in the bracket are simple Gaussian integrals:
\begin{align}
    \int \D{\xi}\; 
    \e^{\frac{im}{2}\frac{\xi^2}{\Delta t}} 
    = 
    \sqrt{\frac{2\pi i\Delta t}{m}} = \mathcal{N}^{-1} \ .
    \label{eq:path:integral:Gaussian:trick:1}
\end{align}
Observe that this factor exactly cancels the annoying normalization in \eqref{eq:path:integral:normalization}---but only for the $\xi$-independent terms in the bracket. The remaining term is a moment of the Gaussian integral. There happens to be a convenient trick for this by taking a derivative of 
\begin{align}
    G_{a,J} \equiv \int \D{x}\; e^{-\frac{a}{2}x^2 + Jx} = \sqrt{\frac{2\pi}{a}}  e^{\frac{J^2}{2a}} 
    \tag{\ref{eq:Gaussian:with:source}} \ ,
\end{align}
where we see that
\begin{align}
    -2\frac{\D{G_{a,0}}}{\D{a}} &= \int \D{x}\; x^2 e^{-\frac{a}{2}x^2} 
    = \frac{1}{a}\sqrt{\frac{2\pi}{a}} \ .
    \label{eq:Gaussian:trick:diff:quad}
\end{align}
\begin{exercise}
While the method above reconstructs even moments of the Gaussian integral, one can reconstruct all moments by taking derivatives with respect to the source, $J$. Show that you can reconstruct the second moment of the Gaussian integral, \eqref{eq:Gaussian:trick:diff:quad}, by taking two derivatives with respect to $J$. 
\end{exercise}
This gives us 
\begin{align}
    \int \D{\xi}\; \xi^2
    \e^{\frac{im}{2}\frac{\xi^2}{\Delta t}} 
    = 
    \frac{1}{-im/\Delta t}
    \sqrt{\frac{2\pi i\Delta t}{m}} 
    =\frac{i\Delta t}{m} \mathcal{N}^{-1} \ .
    \label{eq:path:integral:Gaussian:trick:2}
\end{align}

We can insert \eqref{eq:path:integral:Gaussian:trick:1} and \eqref{eq:path:integral:Gaussian:trick:2} into the expression for $\psi(q,t+\Delta t)$, \eqref{eq:psi:path:integral:schrodinger:intermediate}. The ugly factor of $\mathcal N$ conveniently cancels. We are left with:
\begin{align}
    \psi(q,t+\Delta t) &= \psi(q,t) - i \Delta t V(q)\, \psi(q,t) + \frac{i\Delta t}{2m} \psi''(q,t)
\end{align}
We can rearrange terms and recognize the time derivative:
\begin{align}
    \frac{\D{\psi(q,t)}}{\D{t}}
    &= i\left[\frac{1}{2m}\left(\frac{\textnormal{d}}{\D{q}}\right)^2 - V(q)\right]\psi(q,t) \ ,
\end{align}
which we recognize as the Schr\"odinger equation. You can debate whether or not this counts as a `derivation' of the Schr\"odinger equation since we implicitly assumed the Schr\"odinger equation when writing the time evolution operator. However, if you instead took as a starting point the observation that each path is weighted by $\e^{iS[q(t)]}$, then one would recover the Schr\"odinger equation as a consequence. 








\begin{exercise}
Derive the same results by taking a time derivative of the continuous expression \eqref{eq:path:integral:formalism}. The steps are analogous, but this will test if you understand what the continuous expression means.
\end{exercise}

